\section{Discussion}
\label{sec:discussion}

Research on neuronal selectivity has traditionally focused on canonical cell types---place cells, head direction cells, speed cells--- each characterized with specialized analysis pipelines. While this approach has yielded fundamental insights, it has limited ability to compare selectivity across variable types, to quantify the prevalence of mixed selectivity on a common scale, and to separate genuine encoding from correlation-driven artifacts that arise when behavioral variables covary. These limitations are particularly acute for calcium imaging data, where signals are continuous, noisy, and temporally smeared. Here we address these challenges by systematically screening neuron-behavior pairs spanning spatial, locomotor, and object-related variables within a unified information-theoretic framework. Mutual information quantifies statistical dependence without assuming specific functional forms, enabling INTENSE to accommodate heterogeneous data types---continuous variables like speed and head direction, discrete states like behavioral syllables, and multidimensional quantities like spatial position---within a single analytical framework. 


INTENSE combines Gaussian Copula Mutual Information (GCMI) with circular-shift permutation testing to detect neuron-behavior dependencies in calcium imaging data, achieving balanced performance across a wide range of signal conditions (\fig{fig:synthetics}B,D,E). When applied to the classic benchmark of place cell identification, INTENSE showed high concordance with established event-based methods \cite{skaggs1993information}, converging on an identical ``core'' population of spatially tuned neurons at stringent confidence thresholds (\fig{fig:pc}). The framework operates directly on raw calcium fluorescence without requiring spike deconvolution, avoiding potential distortions or information loss introduced by preprocessing \cite{Rupprecht2021}. The rank-based GCMI estimator manages baseline variations and non-Gaussian transient shapes typical of calcium signal \cite{gcmi}, while two-stage testing enables large-scale analysis by screening out non-significant pairs early. Importantly, synthetic benchmarks underscore a critical requirement for calcium imaging analysis: significance testing requires accounting for temporal autocorrelation (\fig{fig:synthetics}). Without appropriate shuffle controls, the slow dynamics of calcium indicators can inflate apparent selectivity and produce false discoveries \cite{Wei2020}.
Applied to hippocampal activity data during free open field exploration, INTENSE recovers results consistent with key principles of CA1 coding while highlighting the limits of commonly used linear measures. Consistent with sparse coding, only a minority of neurons show significant selectivity to any given behavioral variable, in line with estimates of environment-specific recruitment in CA1 \cite{Wilson1993, Epsztein2011}. Moreover, replacing MI with Pearson correlation in the same pipeline left roughly two-thirds of significant neuron--feature pairs undetected (Section~\ref{sec:mi_vs_corr}), confirming that a substantial fraction of neuron-behavior dependencies are invisible to linear measures. This observation aligns with recent work emphasizing that neural representations occupy curved, nonlinear subspaces, motivating analytical tools that do not assume simple parametric coupling \cite{Altan2021, De2023, intdim}.
A major goal of INTENSE is not only to detect selectivity, but to improve interpretability in the presence of correlated behavioral variables. Temporal autocorrelation and correlated covariates can generate apparent neuron--behavior associations that are statistically real yet misleading with respect to what a neuron encodes \cite{Harris2021nonsense}. More broadly, behavioral variability and uninstructed movements can dominate neural activity and confound selectivity estimates \cite{Stringer2019, Musall2019}. 
INTENSE addresses these concerns through multivariate information analysis that separates true mixed selectivity from apparent selectivity arising from behavioral correlations. Conditional mutual information tests whether selectivity to a variable persists after controlling for correlated behaviors (for example, whether apparent speed tuning remains when controlling for locomotion state or position) \cite{kropff2015speed}. Interaction information further provides a compact summary of whether variables contribute redundantly or synergistically, helping to distinguish apparent mixed selectivity driven by covariance from encoding that reflects conjunction-like computation. This disentanglement becomes particularly important for calcium imaging, where lower temporal resolution already reduces apparent mixed selectivity compared to electrophysiology \cite{Wei2020}.
Applying this disentanglement procedure, INTENSE reliably identifies established selective neurons --- place cells, head direction cells, speed-modulated neurons, object cells --- confirming calcium imaging captures known functional specializations. At the same time, the framework refines estimates of mixed selectivity in naturalistic conditions: approximately a fifth of selective neurons initially appeared multi-selective, but conditional analysis attributed roughly a third of these associations to behavioral redundancy (Section~\ref{sec:intense_deciphers}), consistent % [AUTHOR: find references for mixed coding prevalence]
with distributed mixed coding in the hippocampus. The effect was particularly striking for speed selectivity, where nearly all apparent tuning was explained by locomotion-state correlations. 
Some variables remain fundamentally ambiguous under common experimental constraints: when objects occupy fixed locations, place and object selectivities cannot be clearly separated \cite{Plusnin2021}, and such cases may plausibly reflect genuine conjunctive encoding. More broadly, these results suggest that studies quantifying mixed selectivity in hippocampus, which typically do not control for such behavioral confounds, may overestimate the degree of conjunctive encoding \cite{TyeMiller2024}.
These findings are compatible with theoretical perspectives proposing that apparent complexity can arise from a combination of sparse selectivity at the single-neuron level and structured population codes \cite{Keinath2025}. Although neurons often appear mixed-selective under traditional analysis, our analysis indicates that most encode a single variable. This implies that apparent population-level disentanglement can arise from largely univariate neuronal encoding. The observed sparsity aligns with the log-dynamic brain hypothesis \cite{Buzsaki2014}, where networks organize around minorities of highly active neurons serving as integration points \cite{Fusi2016}. Strikingly, analogous motifs have been discussed in deep artificial networks, where specialized neurons with sparse mixed selectivity emerge spontaneously \cite{Johnston2023, templeton2024scaling}. While the biological and artificial settings are not directly comparable, these parallels motivate the broader hypothesis that sparse mixed selectivity may serve as a computationally efficient strategy for balancing efficiency with representational power.
% [AUTHOR: Consider simplifying this paragraph]
Current limitations reflect broader challenges in large-scale systems neuroscience. Although two-stage testing improves scalability, information-theoretic analyses remain computationally demanding as the number of neurons, behavioral variables, and candidate interactions grows. In addition, our disentanglement procedure can only control for measured covariates; unobserved or poorly captured behavioral factors can still induce residual confounds. The choice of the GCMI estimator represents a deliberate trade-off between nonlinearity detection and computational tractability. GCMI captures monotonic statistical dependencies via rank correlation and enables FFT-accelerated permutation testing and copula reuse across shifts, making large-scale screening feasible. However, non-monotonic tuning curves---such as bell-shaped selectivity centered near the feature median---may be underestimated. Non-parametric estimators such as KSG \cite{Kraskov2004} can capture such dependencies but are substantially more data-demanding and become prohibitively expensive when local non-uniformity correction is applied \cite{Gao2015LNC}, limiting their use in high-throughput pipelines. 
Finally, calcium indicator kinetics and imperfect temporal alignment between neural signals and behavioral variables impose practical constraints on lag selection and temporal resolution, which can influence estimates of selectivity and redundancy. Future integration of Partial Information Decomposition could provide a more granular partition of information into unique, redundant, and synergistic components across multiple variables, strengthening mechanistic interpretation of mixed selectivity beyond pairwise and conditional tests \cite{Williams2010}.
Overall, INTENSE provides an information-theoretic framework tailored to the demands of calcium imaging in continuous free-behavior recordings, aimed at distinguishing genuine selectivity from correlation artifacts. Its generality across neural signal and behavioral variable classes makes it well suited for integration with modern automated phenotyping pipelines. By disentangling behavioral confounds, the method indicates that computational complexity is likely to arise from small, sparse populations of genuinely mixed-selective neurons, rather than from widespread conjunctive encoding across the neural population. As recordings expand to larger populations and increasingly rich behavioral paradigms, correlation-aware selectivity analyses will be increasingly essential for making interpretable claims about neural codes and for avoiding conflation of true encoding with correlation artifacts.
